\section{General conventions and ideas}

This document is, in theory, only a nice-to-have as all relevant information \textit{should} be written in the corresponding style files. However.

\subsection{Introduction}

This collection of macros and classes was started during bachelor programmes (Physics and Information and Computer Engineering) for homework assignments and lab-reports. Based on the commit history, we started in 2020, but it is likely that some forms of the collection were individually started earlier.

The goal of \prettytex{} is to make writing pretty (thus the name) \LaTeX documents easy. This is achieved by outsourcing things like package loading, page layout, colours, macros, operators and such to macro-collections grouped into different packages, which are:

\begin{multicols}{3}
    \begin{itemize}
        \item \nameref{sec:base}
        \item \nameref{sec:math}
        \item gfx
        \item mathematicians
        \item shortcuts
        \item theorems
        \item boxes
        \item scriptum
        \item \warn{invoice}
        \item \highlight{code}
        \item \warn{contract}
        \item \warn{neturalThesis}
        \item \highlight{gfx-components}
        \item \warn{gfx-circuits}
        \item \warn{math-sigtrans}
        \item \error{math-theorems}
        \item \warn{physics}
        \item \warn{pseudo}
        \item \highlight{thesis}
        \item \warn{todo}
    \end{itemize}
\end{multicols}

Packages highlighted in \error{red} are bound to be removed from the main branch, packages marked \warn{orange} are put on the back-burner in the refactor, while a \highlight{purple} highlight means, that the package is going to be refactored. 




\subsection{Package options}

We provide some elementary package options, which are handled via \package{xkeyval}. We have some key-value options, like lang in \nameref{sec:base}, but also quite a few boolean-flags, like vargreek in \nameref{sec:math}. In order to pass a key-value option, load the package as such:

\begin{verbatim}
\usepackage[<option>=<value>]{prettytex/<package>}
\end{verbatim}

To set a boolean-flag to it's non-default value (every boolean flag has a default which is either \verb|true| or \verb|false|), simply pass the option name as such:

\begin{verbatim}
\usepackage[<booleanFlag>]{prettytex/<package>}
\end{verbatim}


\subsection{\prettytex{} macros}

Each package typically provides a few so called \emph{prettytex-macros}. These are used to keep macros, operators and related things modular to a certain degree. This concept is heavily used in \pmath{}. For example, currently, the \verb|\trace| operator results in $\trace$. However, if you wish to alternatively use, say, $\mathrm{tr}$, instead of redefining the \verb|\trace| operator, you can redefine \verb|\prettytexMathTrace| to tr and be done with it. \prettytex{}-macros follow a very simple syntax:

\begin{verbatim}
\prettytex<Package><Macro>,
\end{verbatim}

so all \prettytex{}-macros in \pmath{} are of the form \verb|\prettytexMath<Macro>|. The exception here is \pbase{}, where it simply is \verb|\prettytex<Macro>|. \prettytex{} macros are always listed in their corresponding package description. To keep everything legible, we will mostly omit the \verb|prettytex<Package>| when talking about prettytex-macros, thus we write \prettytex-macros in colour, so \prettytexMacro{Trace} is actually \verb|\prettytexMathTrace|.


\subsection{Convention on argument naming in the documentation}

Many macros use arguments, either mandatory or optional. Most of these can by any valid tex-expression (depending on the mode either text or math), however to highlight their intended use, we name arguments in a certain manner. For example, the argument of a formatting macro will be labelled \verb|text|, while for general mathematical expressions we use \verb|expression|. Optional arguments are highlighted in two ways: either with (o) or (O) (this losely follows the convention of \textsl{xparse}), where (o) means the argument is optional, and (O) optional, but a default value is provided.

\begin{table}[H]
    \centering
    \begin{tabular}{r|p{12cm}}
        \hline\textbf{argument type} & \textbf{description}\\\hline
        \argu{text} & text argument\\
        \argu{colour} & \package{xcolor} colour identifier\\
        \argu{<number type>} & can be \verb|integer| or \verb|float|\\
        \argu{<signed> <number type>} & the sign can be one of \verb|positive| ($>0$), \verb|negative| ($<0$), \verb|nonpositive| $(\leq 0$) or \verb|nonnegative| ($\geq 0$)\\
        \argu{<number type> in <set>} & <set> can be any set, typically used for intervals like \verb|[1,26]|\\
        \argu{expression} & any valid math-mode expression\\
        \argu{array} & valid array-contents (anything that fits into bmatrix and relatives, as well as a plain array in math mode etc)\\
        \argu{reference} & valid reference identifier for labels
    \end{tabular}
    \caption{Argument types}
\end{table}

Some macros allow multiple argument types, these are separated by an or in the argument list, see for example cw in \cref{tab:base_misc}. If a macro uses multiple arguments, they are listed in the expected order and separated by commas.